\documentclass[12pt,a4paper]{article}
\usepackage[T1]{fontenc}
\usepackage[utf8]{inputenc}
\usepackage[spanish]{babel}
\usepackage{geometry}
\usepackage{booktabs}
\usepackage{xcolor}
\usepackage{listings}
\usepackage{hyperref}
\usepackage{enumitem}
\usepackage{titlesec}
\usepackage{parskip}

\geometry{margin=2.5cm}

\hypersetup{
  colorlinks=true,
  linkcolor=black,
  urlcolor=blue!70!black,
}

\lstdefinestyle{bash}{
  language=bash,
  basicstyle=\small\ttfamily,
  backgroundcolor=\color{gray!10},
  frame=single,
  rulecolor=\color{gray!40},
  breaklines=true,
  showstringspaces=false,
}

\lstdefinestyle{c}{
  language=C,
  basicstyle=\small\ttfamily,
  backgroundcolor=\color{gray!10},
  frame=single,
  rulecolor=\color{gray!40},
  breaklines=true,
  showstringspaces=false,
  keywordstyle=\color{blue!70!black},
  commentstyle=\color{green!50!black}\itshape,
  stringstyle=\color{red!60!black},
  morekeywords={u8,__le16,__be16,__le32,__packed,int,static},
}

\title{%
  \textbf{Resolución de problemas de drivers:}\\
  \textbf{NVIDIA y Bluetooth}\\[0.4em]
  \large Compatibilidad de kernel, degradación de driver y\\
  parche DKMS para el ASUS ROG Zephyrus G14 GA403UV
}
\author{pc-config}
\date{Mayo de 2026}

\begin{document}

\maketitle
\tableofcontents
\newpage

% ─────────────────────────────────────────────
\section{Introducción}

Este documento recoge la resolución de dos problemas independientes de drivers
que afectaban al \textbf{ASUS ROG Zephyrus G14 GA403UV} con
\textbf{CachyOS} (base Arch Linux) y el entorno \textbf{Omarchy} (Hyprland +
Wayland):

\begin{enumerate}
  \item \textbf{Driver NVIDIA}: artefactos visuales graves en juegos y
        incompatibilidad del driver 580 con el kernel 7.x, que obligó a fijar
        el kernel LTS~6.18 como arranque predeterminado.
  \item \textbf{Bluetooth MediaTek MT7922}: fallo sistemático en cada arranque
        con el mensaje \texttt{Failed to send wmt func ctrl (-22)}, causado por
        un cambio de comportamiento en el firmware de febrero~de~2026.
\end{enumerate}

Ambos problemas conviven: la solución del Bluetooth requiere el kernel LTS, que
a su vez es necesario por la limitación del driver NVIDIA.

% ─────────────────────────────────────────────
\section{Entorno del sistema}

\begin{center}
\begin{tabular}{ll}
  \toprule
  Componente & Versión / valor \\
  \midrule
  Portátil            & ASUS ROG Zephyrus G14 GA403UV \\
  Sistema operativo   & CachyOS (base Arch Linux) \\
  Entorno de escritorio & Omarchy (Hyprland + Wayland) \\
  Kernel activo       & \texttt{linux-cachyos-lts} 6.18.31-1 \\
  iGPU                & AMD Radeon RX 780M (HawkPoint) --- driver \texttt{amdgpu} \\
  dGPU                & NVIDIA GeForce RTX 4060 Laptop (8\,GB) \\
  Driver NVIDIA       & \texttt{nvidia-open-dkms} 580.95.05-1 \\
  Chip Bluetooth      & MediaTek MT7922 (USB interno) \\
  Firmware BT         & \texttt{linux-firmware} 20260410-1 (build \texttt{20260224103448}) \\
  Gestor de arranque  & Limine 12.2.0 \\
  \bottomrule
\end{tabular}
\end{center}

% ─────────────────────────────────────────────
\section{Problema 1: Driver NVIDIA --- artefactos y compatibilidad de kernel}

\subsection{Síntoma}

Con el driver \texttt{nvidia-open-dkms} \textbf{595.71.05} instalado en el
kernel \texttt{linux-cachyos} (rama~7.x), el juego
\emph{The Last of Us Part~I} y otros títulos mostraban \textbf{artefactos
visuales graves}: corrupciones de textura, píxeles errantes y parpadeos
aleatorios que hacían el juego injugable.

\subsection{Causa}

La versión 595.71.05 de \texttt{nvidia-open} introdujo un bug de renderizado
conocido que afecta a ciertos juegos bajo Wayland/PRIME Offload en
combinación con la iGPU AMD. No se trataba de un problema de configuración
sino de un defecto en el propio driver.

\subsection{Solución: degradar a nvidia-open 580}

Se degradó el driver a la versión \textbf{580.95.05-1}, que no presenta el
bug:

\begin{lstlisting}[style=bash]
# Descargar la version anterior desde la cache de pacman o un mirror
sudo pacman -U nvidia-open-dkms-580.95.05-1-x86_64.pkg.tar.zst \
               nvidia-utils-580.95.05-1-x86_64.pkg.tar.zst \
               lib32-nvidia-utils-580.95.05-1-x86_64.pkg.tar.zst
\end{lstlisting}

Para evitar que pacman actualice el driver automáticamente, añadir al fichero
\texttt{/etc/pacman.conf}:

\begin{lstlisting}[style=bash]
IgnorePkg = nvidia-open-dkms nvidia-utils lib32-nvidia-utils
\end{lstlisting}

\subsection{Por qué el kernel LTS (6.18) y no el kernel 7.x}
\label{sec:kernel-lts}

El driver \texttt{nvidia-open} 580 \textbf{no compila} contra el kernel~7.x
de CachyOS. El módulo DKMS falla durante la fase de \texttt{make} porque la
API interna del kernel cambió entre las ramas~6.18 y~7.x en interfaces que
el driver 580 utiliza directamente (gestión de memoria de vídeo, interfaz
\texttt{drm}, etc.).

El intento de construcción produce errores del tipo:

\begin{lstlisting}[style=bash]
# En /var/lib/dkms/nvidia/580.95.05/build/make.log:
error: implicit declaration of function 'drm_...'
error: unknown type name 'struct drm_...'
\end{lstlisting}

Como resultado, el módulo \texttt{nvidia.ko} no se genera y la dGPU queda
inaccesible. La única versión de driver compatible con la dGPU RTX~4060 que
compila correctamente en CachyOS en el momento de esta guía es la
\textbf{595.x} (que tiene el bug de renderizado) o la \textbf{580.x} (que
requiere el kernel LTS~6.18).

La decisión fue: usar \textbf{nvidia-open 580 + kernel LTS 6.18}.

\subsection{Establecer el kernel LTS como arranque predeterminado}

CachyOS usa \textbf{Limine} como gestor de arranque. Para que
\texttt{linux-cachyos-lts} sea la entrada predeterminada en cada arranque,
editar \texttt{/etc/default/limine} y colocar la entrada LTS antes que
\texttt{linux-cachyos} en \texttt{BOOT\_ORDER}. Luego regenerar la
configuración de Limine con \texttt{sudo limine-update}:

\begin{lstlisting}[style=bash]
sudoedit /etc/default/limine
BOOT_ORDER="linux-cachyos-lts, linux-cachyos, *fallback, Snapshots"
sudo limine-update
\end{lstlisting}


Asegurarse de que la sección del kernel LTS aparece antes que cualquier otra
entrada \texttt{:OS\_ENTRY:}. Limine arranca la primera entrada válida que
encuentre. Verificar tras reiniciar:

\begin{lstlisting}[style=bash]
uname -r
# Resultado esperado: 6.18.31-1-cachyos-lts
\end{lstlisting}

% ─────────────────────────────────────────────
\section{Problema 2: Bluetooth MediaTek MT7922 --- fallo en cada arranque}

\subsection{Síntoma}

Tras cada arranque, el adaptador Bluetooth MT7922 fallaba al inicializarse.
El kernel registraba el siguiente error en \texttt{dmesg}:

\begin{lstlisting}[style=bash]
Bluetooth: hci0: Failed to send wmt func ctrl (-22)
\end{lstlisting}

El código de error \texttt{-22} corresponde a \texttt{-EINVAL}. El Bluetooth
era completamente inoperativo: \texttt{bluetoothctl} no mostraba ningún
adaptador, y no era posible conectar ningún dispositivo.

\subsection{Contexto del driver}

El chip MT7922 es gestionado por los módulos \texttt{btmtk} y \texttt{btusb}
del kernel. El fichero de firmware es:

\begin{lstlisting}[style=bash]
/lib/firmware/mediatek/BT_RAM_CODE_MT7922_1_1_hdr.bin.zst
\end{lstlisting}

La secuencia de arranque del driver es:
\begin{enumerate}
  \item Descarga del firmware al chip (\emph{firmware download}).
  \item Envío del comando \textbf{FUNC\_CTRL enable}: activa el protocolo
        Bluetooth en el firmware.
  \item Verificación del estado: el driver lee la respuesta para confirmar
        que el chip está activo.
\end{enumerate}

El fallo ocurría exactamente en el paso~2.

\subsection{Análisis de la causa raíz}

El protocolo de comunicación entre el driver y el firmware se denomina
\textbf{WMT} (\emph{Wireless Module Testability}). Cada comando WMT espera
una respuesta con una estructura concreta.

Para el comando \texttt{FUNC\_CTRL} (activar Bluetooth), la respuesta
esperada es la estructura \texttt{btmtk\_hci\_wmt\_evt\_funcc}, de
\textbf{9~bytes}:

\begin{center}
\begin{tabular}{lll}
  \toprule
  Campo & Tamaño & Descripción \\
  \midrule
  \texttt{hci\_event\_hdr} & 2\,B & Cabecera HCI (código de evento + longitud) \\
  \texttt{btmtk\_wmt\_hdr} & 5\,B & Cabecera WMT (dir, op, dlen, flag) \\
  \texttt{status}          & 2\,B & Estado del firmware (\texttt{0x0404} = activo) \\
  \bottomrule
\end{tabular}
\end{center}

El firmware con fecha de compilación \textbf{20260224103448} (febrero de
2026, distribuido con \texttt{linux-firmware} 20260410-1) cambió su
comportamiento: en respuesta al comando de activación, devuelve únicamente
\textbf{7~bytes} --- la cabecera completa, pero \emph{sin} los 2~bytes de
estado al final.

El driver intenta extraer esos 2~bytes de estado con \texttt{skb\_pull\_data}:

\begin{lstlisting}[style=c]
case BTMTK_WMT_FUNC_CTRL:
    if (!skb_pull_data(data->evt_skb,
                       sizeof(wmt_evt_funcc->status))) {
        err = -EINVAL;  /* falla: 0 bytes restantes */
        goto err_free_skb;
    }
\end{lstlisting}

Como el \emph{socket buffer} ya está vacío tras haber consumido la cabecera
de 7~bytes, \texttt{skb\_pull\_data} devuelve \texttt{NULL}, la función
retorna \texttt{-EINVAL}, y el driver registra el error.

\textbf{Clave}: el firmware \emph{sí} procesa el comando y activa el
Bluetooth. El cambio es únicamente en el formato de la respuesta. El error
es del driver, no del firmware.

\subsection{La función \texttt{btmtk\_usb\_func\_query}}

El driver ya disponía de una función separada para \emph{consultar} el estado
del Bluetooth sin activarlo, usando el flag~4 del comando FUNC\_CTRL:

\begin{lstlisting}[style=c]
static int btmtk_usb_func_query(struct hci_dev *hdev)
{
    struct btmtk_hci_wmt_params wmt_params;
    int status, err;
    u8 param = 0;

    wmt_params.op   = BTMTK_WMT_FUNC_CTRL;
    wmt_params.flag = 4;   /* consulta, no activacion */
    wmt_params.dlen = sizeof(param);
    wmt_params.data = &param;
    wmt_params.status = &status;

    err = btmtk_usb_hci_wmt_sync(hdev, &wmt_params);
    if (err < 0) {
        bt_dev_err(hdev, "Failed to query function status (%d)", err);
        return err;
    }
    return status; /* BTMTK_WMT_ON_DONE = 5 si el BT esta activo */
}
\end{lstlisting}

Esta consulta con flag~4 devuelve la respuesta completa de 9~bytes
(incluyendo los bytes de estado) incluso con el firmware nuevo, por lo que
funciona correctamente.

\subsection{Solución: módulo DKMS \texttt{btmtk-fix}}

Como el módulo \texttt{btmtk.ko} del kernel no puede modificarse directamente
sin recompilar el kernel completo, se creó un módulo DKMS que reemplaza al
original. Los módulos en \texttt{/updates/dkms/} tienen prioridad automática
sobre los del kernel.

\subsubsection{Estructura del módulo}

\begin{lstlisting}[style=bash]
/usr/src/btmtk-fix-1.0/
    btmtk.c        # copia de drivers/bluetooth/btmtk.c con el parche
    btmtk.h        # cabecera con BTMTK_FIRMWARE_DL_RETRY anadido
    Makefile
    dkms.conf
\end{lstlisting}

\textbf{Makefile}:
\begin{lstlisting}[style=bash]
obj-m := btmtk.o
KDIR ?= /lib/modules/$(shell uname -r)/build
all:
        $(MAKE) -C $(KDIR) M=$(PWD) modules
clean:
        $(MAKE) -C $(KDIR) M=$(PWD) clean
\end{lstlisting}

\textbf{dkms.conf}:
\begin{lstlisting}[style=bash]
PACKAGE_NAME="btmtk-fix"
PACKAGE_VERSION="1.0"
BUILT_MODULE_NAME[0]="btmtk"
DEST_MODULE_LOCATION[0]="/updates/dkms"
AUTOINSTALL="yes"
\end{lstlisting}

\subsubsection{El parche aplicado}

En \texttt{btmtk.h}, se añadió la constante de flag necesaria para el
mecanismo de reintento:

\begin{lstlisting}[style=c]
enum {
    /* ... constantes existentes ... */
    BTMTK_ISOPKT_RUNNING,
    BTMTK_FIRMWARE_DL_RETRY,   /* anadido: indica reintento en curso */
};
\end{lstlisting}

En \texttt{btmtk.c}, dentro de \texttt{btmtk\_usb\_setup}, tras el envío del
comando de activación para los chips MT7922/7925/7961, se añadió el bloque
de recuperación cuando la respuesta es \texttt{-EINVAL}:

\begin{lstlisting}[style=c]
err = btmtk_usb_hci_wmt_sync(hdev, &wmt_params);
if (err == -EINVAL) {
    /* El firmware nuevo (feb 2026+) no incluye bytes de estado en
     * la respuesta FUNC_CTRL enable. En su lugar, consultar el
     * estado via polling, tolerando -EINVAL mientras el firmware
     * termina de inicializarse tras una carga en frio (hasta ~2s).
     */
    int qstatus;

    err = readx_poll_timeout(btmtk_usb_func_query, hdev, qstatus,
                             qstatus != -EINVAL &&
                             qstatus != BTMTK_WMT_ON_PROGRESS,
                             100000, 3000000);
    if (!err) {
        if (qstatus == BTMTK_WMT_ON_DONE)
            err = 0;
        else
            err = (qstatus < 0) ? qstatus : -EIO;
    }
    bt_dev_info(hdev, "FUNC_CTRL status poll: err=%d qstatus=%d",
                err, qstatus);
}
if (err < 0) {
    if (!test_and_set_bit(BTMTK_FIRMWARE_DL_RETRY, &btmtk_data->flags))
        btmtk_reset_sync(hdev);
    bt_dev_err(hdev, "Failed to send wmt func ctrl (%d)", err);
    return err;
}
\end{lstlisting}

La macro \texttt{readx\_poll\_timeout} (de \texttt{<linux/iopoll.h>}) llama
a \texttt{btmtk\_usb\_func\_query} repetidamente cada 100\,ms durante un
máximo de 3\,s. La condición de salida es que la respuesta no sea
\texttt{-EINVAL} (firmware aún inicializándose) ni
\texttt{BTMTK\_WMT\_ON\_PROGRESS} (carga en curso). Cuando la respuesta es
\texttt{BTMTK\_WMT\_ON\_DONE}~(=5), el Bluetooth está activo y el setup
continúa con éxito.

\subsubsection{Comandos de construcción e instalación}

\begin{lstlisting}[style=bash]
# Registrar el modulo en DKMS
sudo dkms add /usr/src/btmtk-fix-1.0

# Compilar para el kernel activo
sudo dkms build btmtk-fix/1.0 -k 6.18.31-1-cachyos-lts

# Instalar (copia el .ko.zst a /updates/dkms/, tiene prioridad)
sudo dkms install btmtk-fix/1.0 -k 6.18.31-1-cachyos-lts
\end{lstlisting}

Para reconstruir desde cero si es necesario:

\begin{lstlisting}[style=bash]
sudo dkms remove btmtk-fix/1.0 --all
sudo dkms add /usr/src/btmtk-fix-1.0
sudo dkms build  btmtk-fix/1.0 -k 6.18.31-1-cachyos-lts
sudo dkms install btmtk-fix/1.0 -k 6.18.31-1-cachyos-lts
\end{lstlisting}

El módulo queda instalado en:
\begin{lstlisting}[style=bash]
/lib/modules/6.18.31-1-cachyos-lts/updates/dkms/btmtk.ko.zst
\end{lstlisting}

El módulo original del kernel queda en:
\begin{lstlisting}[style=bash]
/lib/modules/6.18.31-1-cachyos-lts/kernel/drivers/bluetooth/btmtk.ko.zst
\end{lstlisting}

\subsection{Verificación}

Recargar el módulo y comprobar \texttt{dmesg}:

\begin{lstlisting}[style=bash]
sudo modprobe -r btusb btmtk
sudo dmesg -C
sudo modprobe btusb
sleep 4
sudo dmesg | grep -iE "hci0.*(HW.SW|FUNC_CTRL|poll|device setup|AOSP|failed)"
\end{lstlisting}

Salida esperada (fallo \textbf{ausente}, setup correcto):

\begin{lstlisting}[style=bash]
Bluetooth: hci0: HW/SW Version: 0x008a008a, Build Time: 20260224103448
Bluetooth: hci0: FUNC_CTRL status poll: err=0 qstatus=5
Bluetooth: hci0: Device setup in 125819 usecs
Bluetooth: hci0: AOSP extensions version v1.00
Bluetooth: hci0: AOSP quality report is supported
\end{lstlisting}

La línea \texttt{FUNC\_CTRL status poll: err=0 qstatus=5} confirma que el
parche está activo y que el polling detectó correctamente el estado
\texttt{ON\_DONE} del firmware.

% ─────────────────────────────────────────────
\section{Notas de mantenimiento}

\subsection{Si se actualiza el kernel LTS}

DKMS recompila automáticamente el módulo \texttt{btmtk-fix} para cada nuevo
kernel gracias a \texttt{AUTOINSTALL="yes"} en \texttt{dkms.conf}. Si la
compilación falla por cambios de API en el kernel, el log de error se
encuentra en:

\begin{lstlisting}[style=bash]
/var/lib/dkms/btmtk-fix/1.0/build/make.log
\end{lstlisting}

\subsection{Si se actualiza linux-firmware}

El parche es compatible con versiones futuras del firmware. Si el firmware
vuelve a incluir los bytes de estado en la respuesta FUNC\_CTRL, \texttt{err}
valdrá 0 (no \texttt{-EINVAL}) y el bloque de polling no se ejecutará. El
comportamiento es retrocompatible.

\subsection{Si se actualiza nvidia-open-dkms}

Antes de actualizar el driver NVIDIA, verificar que la nueva versión compila
contra el kernel LTS~6.18. Consultar el \emph{changelog} del paquete en AUR
o en los foros de CachyOS. Si la nueva versión exige el kernel 7.x, mantener
el \texttt{IgnorePkg} en \texttt{pacman.conf} hasta disponer de un driver
compatible.

\subsection{Verificar el estado en cada arranque}

Para confirmar que ambos componentes funcionan tras un reinicio:

\begin{lstlisting}[style=bash]
# Kernel activo (debe ser 6.18.x-cachyos-lts)
uname -r

# Modulo NVIDIA cargado
lsmod | grep nvidia

# Bluetooth activo (debe mostrar el adaptador)
bluetoothctl show | grep -E "Name|Powered|Discovering"

# Sin errores de BT en dmesg
sudo dmesg | grep -i "bluetooth.*failed"
\end{lstlisting}

\end{document}
